\documentclass[10pt,usepdftitle=false,aspectratio=169]{beamer}

%!TEX root = talk.tex

% serif fonts for mathematics
%\usefonttheme[onlymath]{serif}

\usepackage{animate}
\usepackage{textcase,marvosym,cancel,booktabs,colortbl,nicefrac,ifsym}
\usepackage{fourier-orns}
\usepackage{amsmath,amssymb,graphicx,tikz,pgfplots,colonequals,mathtools} % MnSymbol
\usepackage{booktabs,multirow,multicol}
\usepackage{pifont} % for the dingbats in the line search plots
% \usepackage{stix} % for \sumint

\usepackage{multido}
\usepackage{makecell}
\usepackage{soul} % for better underline with ul{}
\usetikzlibrary{tikzmark,decorations.pathreplacing,calc}

%% not used with xetex
\usepackage{microtype}
\usepackage[utf8]{inputenc}

% a few colours, for general use.
\definecolor{lred}{RGB}{200,0,0}
\definecolor{dred}{RGB}{130,0,0} 
\definecolor{dblu}{RGB}{0,0,130}
\definecolor{dgre}{RGB}{0,130,0} 
\definecolor{dgra}{RGB}{50,50,50}
\definecolor{mgra}{RGB}{221,222,214}
\definecolor{lgra}{RGB}{238,238,234}
\definecolor{MPG}{RGB}{000,125,122}
\definecolor{lMPG}{RGB}{000,190,189}
\definecolor{ora}{HTML}{cc7a00} 
\definecolor{lblu}{HTML}{7DA7D9}%PS blue

% Color scheme of the Eberhard-Karls University
\definecolor{TUred}{RGB}{141,45,57}
\definecolor{TUdark}{RGB}{55,65,74}
\definecolor{TUgold}{RGB}{174,159,109}
\definecolor{TUgray}{RGB}{175,179,183}

\definecolor{ERC_ora}{RGB}{233,93,15}

\setlength{\parindent}{0pt}

\newcommand{\mpg}[1]{{\color{MPG} #1}}   % highlight command 1
\newcommand{\dre}[1]{{\color{TUred} #1}}   % highlight command 1
\newcommand{\blu}[1]{{\color{dblu} #1}}   % highlight command 1
\newcommand{\ora}[1]{{\color{ora} #1}}   % highlight command 1
\newcommand{\gra}[1]{{\color{mgra} #1}}   % highlight command 1
\newcommand{\gold}[1]{{\color{TUgold} #1}}   % highlight command 1
\setbeamercolor{alerted text}{fg = TUred} % highlight command 2

\setbeamercolor{normal text}{fg=black,bg=white}
\setbeamercolor{structure}{fg=TUred}

\setbeamercolor{item projected}{use=item,fg=black,bg = TUred}

\setbeamercolor*{palette primary}{fg=white,bg=TUred}
\setbeamercolor*{palette secondary}{parent=palette primary,use=palette
primary,bg=dblu} 
\setbeamercolor*{palette tertiary}{parent=palette
primary,use=palette primary,fg=white,bg=dgre} 
\setbeamercolor*{palette
quaternary}{parent=palette primary,use=palette primary,bg=dgre}

\setbeamercolor*{block body}{bg=TUgray, fg =black}
\setbeamercolor*{block title}{parent=structure,bg=TUdark,fg=white}
\setbeamercolor{block body alerted}{bg=TUgray,fg=TUred}
\setbeamercolor{block title example}{fg=MPG,bg=white}
\setbeamercolor{block body example}{fg=black,bg=white}

\setbeamercolor{frametitle}{bg=TUdark,fg=white}
\setbeamercolor{frametitle right}{bg=white}
\setbeamercolor{framesubtitle}{fg=TUgray}

% \setbeamerfont{framesubtitle}{size*={8}{12}}

\setbeamercolor{title}{fg=TUred}
\setbeamercolor{subtitle}{fg=black} \setbeamercolor{author}{fg=TUdark}
\setbeamercolor{date}{fg=TUdark}

\setbeamercolor*{titlelike}{parent=structure}

\setbeamertemplate{navigation symbols}{}
\setbeamertemplate{bibliography item}[triangle]

\setbeamerfont{frametitle}{}

%%%% REDEFINE \emph to COLORS.
\let\emph\relax % there's no \RedeclareTextFontCommand
\DeclareTextFontCommand{\emph}{\color{TUred}\bfseries} % changed from \em

% \setbeamertemplate{itemize items}{\color{TUgold}$\blacktriangleright$}
\setbeamertemplate{itemize items}{\starredbullet}

\setbeamertemplate{footline}{\hfill\color{TUdark}{\insertframenumber}\hspace{2ex}\null\newline\vspace{2mm}}

\arrayrulecolor{TUdark}

\newcommand{\filltotal}{\hspace{0pt plus 1 filll}}

\newcommand{\framesubtitlecite}[2]{\framesubtitle{#1 \filltotal #2}}

\newcommand{\titlemark}[1]{
  \begin{tikzpicture}[remember picture, overlay]
    \node[draw=none,text=TUgray,anchor=north east,yshift=-.8259cm] at (current
    page.north east) {\footnotesize{#1}};
  \end{tikzpicture}
}

\newcommand{\graybox}[1]{
  \tikzexternaldisable%
  \begin{center}%
    \tikz{\node[fill=TUdark,text width=\textwidth]{%
        \begin{minipage}{1.0\linewidth}%
        \color{white}
          \begin{center}%
            #1
          \end{center}%
        \end{minipage}};}%
  \end{center}%
  \tikzexternalenable 
}

\newcommand{\divider}{\noindent\makebox[\linewidth]{\rule{\paperwidth}{.4pt}}  }

\newcommand{\paperwhite}{\includegraphics[height=.6\baselineskip]{../text-document_white.png}\hspace{.5em}}
\newcommand{\paperblack}{\includegraphics[height=.6\baselineskip]{../text-document.png}\hspace{.5em}}

\newcommand{\bookwhite}{\includegraphics[height=.6\baselineskip]{../book-white.png}\hspace{.5em}}
\newcommand{\bookblack}{\includegraphics[height=.6\baselineskip]{../book-black.png}\hspace{.5em}}

\setbeamercolor{ribboncolor}{fg=black,bg=TUgold}
\newcommand{\ribbon}[1]{
    \begin{beamercolorbox}[wd=\paperwidth,colsep*=.3em,center]{ribboncolor}
    \setbeamertemplate{itemize items}{\color{black}\starredbullet}
    \setbeamercolor{structure}{fg=white}
    \begin{minipage}{1.0\textwidth}%
            #1
        \end{minipage}
    \end{beamercolorbox}
}
\setbeamercolor{whiteribboncolor}{fg=black,bg=white}
\newcommand{\whiteribbon}[1]{
    \begin{beamercolorbox}[wd=\paperwidth,colsep*=.5em,center,text width=\textwidth]{whiteribboncolor}
    \setbeamertemplate{itemize items}{\color{black}\starredbullet}
    \setbeamercolor{structure}{fg=black}
    \begin{minipage}{\textwidth}%
            #1
        \end{minipage}
    \end{beamercolorbox}
}

\newcommand{\bfa}[1]{
  \begin{beamercolorbox}[wd=\paperwidth,colsep*=.5em,center]{white}
    \setbeamertemplate{itemize items}{\color{black}\starredbullet}
    \tikzexternaldisable
    \begin{tikzpicture}
    \node[signal,minimum width=\paperwidth,draw=TUgold,fill=TUgold,text=black,text width=\textwidth]{
    \begin{minipage}{\textwidth}%
            #1
        \end{minipage}};
    \end{tikzpicture}
    \tikzexternalenable
    \end{beamercolorbox}
}

\newcommand{\bfi}[1]{
  \begin{beamercolorbox}[wd=\paperwidth,colsep*=.5em,center,text width=\textwidth]{white}
    \setbeamertemplate{itemize items}{\color{black}\starredbullet}
    \tikzexternaldisable
    \begin{tikzpicture}
    \node[signal,signal from=west,signal to=nowhere,minimum width=\linewidth,draw=TUgold,fill=TUgold,text=black,signal pointer angle=140]{
    \begin{minipage}{\textwidth}%
            #1
        \end{minipage}};
    \end{tikzpicture}\hspace{2.5mm}
    \tikzexternalenable
    \end{beamercolorbox}
}

\newcommand{\blackslide}{
{\setbeamercolor{background canvas}{bg=black}
  \begin{frame}[plain]
    \null
  \end{frame}
}}


\newcommand{\blackslidetext}[1]{
{\setbeamercolor{background canvas}{bg=TUdark}%
  \setbeamercolor{structure}{fg=white}%
  \setbeamercolor{normal text}{fg=white}%
  \setbeamercolor{body}{fg=white}%
  \setbeamertemplate{itemize items}{\color{white}\starredbullet}%
  \begin{frame}
    \color{white} #1
  \end{frame}
}}

% logo in the upper right corner
\usepackage{eso-pic}
\newcommand\AtPagemyUpperLeft[1]{\AtPageLowerLeft{%
\put(\LenToUnit{0.83\paperwidth},\LenToUnit{0.915\paperheight}){#1}}}
\AddToShipoutPictureFG{
  \AtPagemyUpperLeft{{\includegraphics[width=2.5cm,keepaspectratio]{assets/UT_WBMW_Weiss_1C.pdf}}}
}%

%%%% REDEFINE \emph to COLORS.
\let\emph\relax % there's no \RedeclareTextFontCommand
\DeclareTextFontCommand{\emph}{\color{TUred}\bfseries} % } changed from \em

%% Macros
\newcommand{\g}{\mid} 
\newcommand{\de}{\partial}
\renewcommand{\d}{\:\mathrm{d}} 
\newcommand{\Exp}{\mathsf{E}}
\newcommand{\eps}{\varepsilon}
\newcommand{\Ent}{\mathsf{H}}
\newcommand{\Cov}{\operatorname{cov}} 
\newcommand{\cov}{\operatorname{cov}} 
\newcommand{\std}{\operatorname{std}} 
\newcommand{\Var}{\operatorname{var}} 
\newcommand{\var}{\operatorname{var}} 
\renewcommand{\H}{\mathcal{H}}
\newcommand{\K}{\mathcal{K}} 
\newcommand{\KL}{\text{KL}} 
\renewcommand{\Re}{\mathbb{R}}
\newcommand{\N}{\mathcal{N}} 
\newcommand{\M}{\mathcal{M}}
\newcommand{\cC}{\mathcal{C}}
\newcommand{\D}{\mathcal{D}}
\newcommand{\B}{\mathcal{B}}
\renewcommand{\L}{\mathcal{L}}
\newcommand{\Trans}{^{\intercal}} 
\newcommand{\rep}{\text{rep}}
\newcommand{\base}{\text{base}} 
\newcommand{\eval}{\text{eval}}
\newcommand{\diag}{\operatorname{diag}}
\newcommand{\argmin}{\operatorname*{arg\:min}}
\newcommand{\argmax}{\operatorname*{arg\:max}}
%\renewcommand{\=}{\operatorname*{=}}
\newcommand{\spa}{\operatorname{span}} 
\newcommand{\erf}{\operatorname{erf}}
\newcommand{\ce}{\colonequals}
\newcommand{\ec}{\equalscolon}
\newcommand{\tr}{\operatorname{tr}}
\newcommand{\rk}{\operatorname{rk}}

\makeatletter
% Taken from http://ctan.org/pkg/centernot
\newcommand*{\centernot}{%
  \mathpalette\@centernot
}
\def\@centernot#1#2{%
  \mathrel{%
    \rlap{%
      \settowidth\dimen@{$\m@th#1{#2}$}%
      \kern.5\dimen@
      \settowidth\dimen@{$\m@th#1=$}%
      \kern-.5\dimen@
      $\m@th#1\not$%
    }%
    {#2}%
  }%
}
\makeatother

\newcommand{\independent}{\perp\mkern-9.5mu\perp}
\newcommand{\notindependent}{\centernot{\independent}}

\newcommand{\q}{\quad}
\newcommand{\qq}{\qquad}
\newcommand{\qqq}{\quad\qquad}
\newcommand{\qqqq}{\qquad\qquad}

\usepackage[e]{esvect} % for extensible vector accents (\vect{})

\renewcommand{\vec}{\boldsymbol} 
\newcommand{\fun}[1]{\mathsf{#1}}
\newcommand{\vect}[1]{\vv{#1}}
% \newcommand{\vect}[1]{\overrightharpoon{#1}}
\renewcommand{\O}{\mathcal{O}}
\newcommand{\GP}{\mathcal{GP}}
\newcommand{\Id}{I} 
\newcommand{\II}{\mathbb{I}}

\makeatletter
\newcommand{\superimpose}[2]{%s
  {\ooalign{$#1\@firstoftwo#2$\cr\hfil$#1\@secondoftwo#2$\hfil\cr}}}
\makeatother
\newcommand{\ostimes}{\mathpalette\superimpose{{\otimes}{\ominus}}}
\newcommand{\oatimes}{\mathpalette\superimpose{{\otimes}{\obar}}}

\DeclareMathOperator*{\sumint}{\mathpalette\superimpose{{\int}{\sum}}}
% \newcommand{\ostimes}{\circledast}

\newcommand{\y}{\vec{y}}
\newcommand{\x}{\vec{x}}
\newcommand{\X}{\vec{X}}
\newcommand{\bX}{\mathbb{X}}
\newcommand{\const}{\mathrm{const.}}

\newcommand{\up}{\textasciicircum}


%%%%%%%% TiKZ %%%%%%%%
\usetikzlibrary{arrows,shapes,plotmarks,decorations.pathmorphing}
\usetikzlibrary{backgrounds,calc,positioning,fadings}

\tikzset{>=stealth'} 
\tikzstyle{graphnode} = 
   [circle,draw=TUdark,minimum size=22pt,text centered,text
     width=22pt,inner sep=0pt] 
\tikzstyle{var}   =[graphnode,fill=white]
\tikzstyle{obs}   =[graphnode,fill=TUdark,text=white]
\tikzstyle{act}   =[rectangle,draw=black,text=white,minimum
size=22pt,text centered, text width=22pt,inner sep=0pt]
\tikzstyle{fac}   =[rectangle,draw=black,fill=black!25,minimum size=5pt]
\tikzstyle{facprior} =[rectangle,draw=black,fill=black,text=white,minimum size=5pt]
\tikzstyle{edge}  =[draw=TUdark,thin,->]
\tikzstyle{undirected}  =[draw=TUdark,thin,-]
\tikzstyle{prior} =[rectangle, draw=black, fill=black, minimum size=
5pt, inner sep=0pt]
\tikzstyle{dirprior} = [circle, draw=black, fill=black, minimum
size=5pt, inner sep=0pt]

\tikzfading[name=fade top,bottom color=transparent!0,top color=transparent!75]

% to avoid warnings, copy only two symbols from stmaryrd
\DeclareSymbolFont{stmry}{U}{stmry}{m}{n}
\DeclareMathSymbol\leftarrowtriangle\mathrel{stmry}{"5E}
\DeclareMathSymbol\rightarrowtriangle\mathrel{stmry}{"5F}
\DeclareMathSymbol\sslash\mathrel{stmry}{"0C}
\DeclareMathSymbol\obar\mathrel{stmry}{"3A}
\DeclareMathSymbol\otimes\mathrel{stmry}{"0F}
\DeclareMathSymbol\ominus\mathrel{stmry}{"17}
\DeclareMathSymbol\minuso\mathrel{stmry}{"0A}
\renewcommand{\gets}{\operatorname*{\leftarrowtriangle}}
\renewcommand{\to}{\operatorname*{\rightarrowtriangle}}

\usetikzlibrary{arrows,shapes,plotmarks,pgfplots.colormaps}
\usetikzlibrary{pgfplots.groupplots}
\pgfplotsset{compat=newest}
\pgfplotsset{
  every axis legend/.append style =
    {
      cells = { anchor = east },
      draw  = none
    },
}  
\makeatletter
\pgfplotsset{ 
    range frame/.style={
        tick align = outside,
        axis line style={opacity=0},
        after end axis/.code={
            \draw ({rel axis cs:0,0}-|{axis cs:\pgfplots@data@xmin,0}) -- ({rel axis cs:0,0}-|{axis cs:\pgfplots@data@xmax,0});
            \draw ({rel axis cs:0,0}|-{axis cs:0,\pgfplots@data@ymin}) -- ({rel axis cs:0,0}|-{axis cs:0,\pgfplots@data@ymax});
        }
    }
}
\makeatother

\pgfkeys{/pgfplots/mystyle/.style={
  % semithick,
  % tick style={major tick length=4pt,semithick,gray},
  xtick align = inside,
  ytick align = inside
  }}

\pgfkeys{/pgfplots/mytuftestyle/.style={
  semithick,
  tick style={major tick length=4pt,semithick,black},
  separate axis lines,
  axis x line*=bottom,
  axis x line shift=5pt,
  xlabel shift=0pt,
  axis y line*=left,
  tick align = outside,
  axis y line shift=5pt,
  ylabel shift=0pt}}


%%%% THEOREM environments
\newtheorem{defi}{Definition}[section]
\newtheorem{theo}[defi]{Theorem}
\newtheorem{cor}[defi]{Corollary}
\newtheorem{lemm}[defi]{Lemma}
\newtheorem{prob}[defi]{Problem}

\newlength{\figureheight}
\newlength{\figurewidth}
\newlength{\figheight}
\newlength{\figwidth}

% algorithms
\usepackage{algorithm}
\usepackage{algpseudocode}
\algrenewcommand{\algorithmiccomment}[1]{\hfill {$\sslash$ \scriptsize #1}}
\algrenewcommand\alglinenumber[1]{\tiny #1}


% %%%%%% adding vertical bars to algorithmicx. Based on a response at 
% % http://tex.stackexchange.com/questions/52473/is-it-possible-to-have-connecting-loop-lines-like-algorithm2e-in-algorithmic/52778#52778
\makeatletter
% This is the vertical rule that is inserted
\def\therule{\makebox[\algorithmicindent][l]{\hspace*{.5em}\vrule height .75\baselineskip depth .25\baselineskip}}%

\newtoks\therules% Contains rules
\therules={}% Start with empty token list
\def\appendto#1#2{\expandafter#1\expandafter{\the#1#2}}% Append to token list
\def\gobblefirst#1{% Remove (first) from token list
  #1\expandafter\expandafter\expandafter{\expandafter\@gobble\the#1}}%
\def\LState{\State\unskip\the\therules}% New line-state
\def\pushindent{\appendto\therules\therule}%
\def\popindent{\gobblefirst\therules}%
\def\printindent{\unskip\the\therules}%
\def\printandpush{\printindent\pushindent}%
\def\popandprint{\popindent\printindent}%

%      ***      DECLARED LOOPS      ***
% (from algpseudocode.sty)
\algdef{SE}[WHILE]{While}{EndWhile}[1]
  {\printandpush\algorithmicwhile\ #1\ \algorithmicdo}
  {\popandprint\algorithmicend\ \algorithmicwhile}%
\algdef{SE}[FOR]{For}{EndFor}[1]
  {\printandpush\algorithmicfor\ #1\ \algorithmicdo}
  {\popandprint\algorithmicend\ \algorithmicfor}%
\algdef{S}[FOR]{ForAll}[1]
  {\printindent\algorithmicforall\ #1\ \algorithmicdo}%
\algdef{SE}[LOOP]{Loop}{EndLoop}
  {\printandpush\algorithmicloop}
  {\popandprint\algorithmicend\ \algorithmicloop}%
\algdef{SE}[REPEAT]{Repeat}{Until}
  {\printandpush\algorithmicrepeat}[1]
  {\popandprint\algorithmicuntil\ #1}%
\algdef{SE}[IF]{If}{EndIf}[1]
  {\printandpush\algorithmicif\ #1\ \algorithmicthen}
  {\popandprint\algorithmicend\ \algorithmicif}%
\algdef{C}[IF]{IF}{ElsIf}[1]
  {\popandprint\pushindent\algorithmicelse\ \algorithmicif\ #1\ \algorithmicthen}%
\algdef{Ce}[ELSE]{IF}{Else}{EndIf}
  {\popandprint\pushindent\algorithmicelse}%
\algdef{SE}[PROCEDURE]{Procedure}{EndProcedure}[2]
   {\printandpush\algorithmicprocedure\ \textproc{#1}\ifthenelse{\equal{#2}{}}{}{(#2)}}%
   {\popandprint\algorithmicend\ \algorithmicprocedure}%
\algdef{SE}[FUNCTION]{Function}{EndFunction}[2]
   {\printandpush\algorithmicfunction\ \textproc{#1}\ifthenelse{\equal{#2}{}}{}{(#2)}}%
   {\popandprint\algorithmicend\ \algorithmicfunction}%
\makeatother

%%%% fiddling with lists:
% \usepackage{enumitem}
% try: \begin{description}[leftmargin=\parindent,labelindent=\parindent] ...

%%%% FONTS

\usepackage{inconsolata}
% \usepackage[defaultsans]{droidsans}
\usepackage[sfdefault,light,condensed]{roboto}
% \renewcommand*\ttdefault{cmvtt} % use with roboto condensed
% \usepackage[sfdefault]{universalis}
% \usepackage[default]{lato}
% \usepackage{lmodern}
% \usepackage{paratype}
% \usepackage{cmbright}
% \usepackage[default,osfigures,scale=0.95]{opensans}
\usepackage{sfmath}
\renewcommand*\familydefault{\sfdefault}
\usepackage[T1]{fontenc}

% special \vneq that works with sfmath
\makeatletter
\newcommand*{\vneq}{%
  \mathrel{%
    \mathpalette\@vneq{=}%
  }%
}
\newcommand*{\@vneq}[2]{%
  % #1: math style (\displaystyle, \textstyle, ...)
  % #2: symbol (=, ...)
  \sbox0{\raisebox{\depth}{$#1\neq$}}%
  \sbox2{\raisebox{\depth}{$#1|\m@th$}}%
  \ifdim\ht2>\ht0 %
    \sbox2{\resizebox{\vneqxscale\width}{\vneqyscale\ht0}{\unhbox2}}%
  \fi
  \sbox2{$\m@th#1\vcenter{\copy2}$}%
  \ooalign{%
    \hfil\phantom{\copy2}\hfil\cr
    \hfil$#1#2\m@th$\hfil\cr
    \hfil\copy2\hfil\cr
  }%
}
\newcommand*{\vneqxscale}{1}
\newcommand*{\vneqyscale}{.8}
\makeatother
\newcommand*{\eg}{e.g.\@\xspace}
\newcommand*{\ie}{i.e.\@\xspace}
\newcommand*{\cf}{cf.\@\xspace}
\newcommand{\amsbound}{\textsc{AMSBound}\xspace}
\newcommand{\adabound}{\textsc{AdaBound}\xspace}
\newcommand{\adam}{\textsc{Adam}\xspace}
\newcommand{\momentum}{\textsc{Momentum}\xspace}
\newcommand{\sgd}{\textsc{SGD}\xspace}
\newcommand{\rmsprop}{\textsc{RMSProp}\xspace}
\newcommand{\amsgrad}{\textsc{AMSGrad}\xspace}
\newcommand{\adadelta}{\textsc{Adadelta}\xspace}
\newcommand{\nag}{\textsc{NAG}\xspace}
\newcommand{\nadam}{\textsc{Nadam}\xspace}
\newcommand{\radam}{\textsc{Radam}\xspace}
\newcommand{\norm}[1]{\left\lVert#1\right\rVert}
\newcommand{\plus}{\raisebox{.4\height}{\scalebox{.6}{+}}}
\newcommand{\minus}{\raisebox{.4\height}{\scalebox{.8}{-}}}
\newcommand{\imagequadsize}{1cm}
\newcommand{\domainsize}[1]{\scriptsize{#1}}
\newcommand{\domainbed}{\textsc{DomainBed}\xspace}
\newcommand{\lime}{\textsc{LIME}}

\newcommand{\featureg}{\mathbf{g}_\mathbf{z}}
\newcommand{\featuregl}{\mathbf{g}_{\mathbf{z},i,j}}
\newcommand{\featurega}{\Tilde{\mathbf{g}}_\mathbf{z}}
\newcommand{\featuregal}{\Tilde{\mathbf{g}}_{\mathbf{z}, i, j}}
\newcommand{\featuregac}{\Tilde{\mathbf{g}}_{\mathbf{z},c}}
\newcommand{\featuregcl}{\mathbf{g}_{\mathbf{z},c,i,j}}
\newcommand{\loguni}[2]{\mathcal{L}\mathcal{U}_{10}(#1, #2)}
\newcommand{\logunitwo}[2]{\mathcal{L}\mathcal{U}_{2}(#1, #2)}
\newcommand{\logunix}[2]{\mathcal{L}\mathcal{U}_{x}(#1, #2)}
\newcommand{\uni}[2]{\mathcal{U}(#1, #2)}

\newcommand{\rsc}{\textsc{RSC}\xspace}
\newcommand{\scam}{\textsc{SCAM}\xspace}
\newcommand{\scamb}{\textsc{SCAM-B}\xspace}
\newcommand{\scambs}{\textsc{SCAM-BS}\xspace}
\newcommand{\scamdb}{\textsc{SCAM-DB}\xspace}
\newcommand{\scamdbs}{\textsc{SCAM-DBS}\xspace}
\newcommand{\scamc}{\textsc{SCAM-C}\xspace}
\newcommand{\scamcs}{\textsc{SCAM-CS}\xspace}
\newcommand{\scamdc}{\textsc{SCAM-DC}\xspace}
\newcommand{\scamdcs}{\textsc{SCAM-DCS}\xspace}
\newcommand{\scamt}{\textsc{SCAM-T}\xspace}
\newcommand{\scamts}{\textsc{SCAM-TS}\xspace}
\newcommand{\scamdt}{\textsc{SCAM-DT}\xspace}
\newcommand{\scamdts}{\textsc{SCAM-DTS}\xspace}
\newcommand{\tscam}{\textsc{SCAM*}\xspace}
\newcommand{\tscamb}{\textsc{SCAM-B*}\xspace}
\newcommand{\tscambs}{\textsc{SCAM-BS*}\xspace}
\newcommand{\tscamdb}{\textsc{SCAM-DB*}\xspace}
\newcommand{\tscamdbs}{\textsc{SCAM-DBS*}\xspace}
\newcommand{\tscamc}{\textsc{SCAM-C*}\xspace}
\newcommand{\tscamcs}{\textsc{SCAM-CS*}\xspace}
\newcommand{\tscamdc}{\textsc{SCAM-DC*}\xspace}
\newcommand{\tscamdcs}{\textsc{SCAM-DCS*}\xspace}
\newcommand{\tscamt}{\textsc{SCAM-T*}\xspace}
\newcommand{\tscamts}{\textsc{SCAM-TS*}\xspace}
\newcommand{\tscamdt}{\textsc{SCAM-DT*}\xspace}
\newcommand{\tscamdts}{\textsc{SCAM-DTS*}\xspace}

\newcommand{\tabtop}[1]{\textcolor{darkred}{\mathbf{#1}}}

\DeclarePairedDelimiter\ceil{\lceil}{\rceil}
\DeclarePairedDelimiter\floor{\lfloor}{\rfloor}

\usetikzlibrary{external}
%\tikzexternalize[mode=list and make]
% use  make -j 8 -f talk.makefile to compile with 8 parallel threads (this is what it takes to max out the machine, depsite it having 4 cores)
%\tikzset{external/force remake=false}
%\tikzsetexternalprefix{fig/external/}

\begin{document}

\tikzexternaldisable
\begin{frame}
  \title{{\bf Explainability-aided Domain Generalization}\newline Master Thesis
  \vspace*{-.7cm}}
  \author{Robin M. Schmidt \vspace{-1cm}} \date{02 March 2021}

  \vspace{-1.5cm}
  \maketitle 
  \vspace{-1.0cm}

  \begin{center}
      \includegraphics[width=0.4\textwidth]{assets/UT_WBMW_Rot_RGB.pdf}
  \end{center}

  \thispagestyle{empty}
  \setcounter{framenumber}{0}

  %%%%%%%%%%%%%%%% ANIMATED LOGO %%%%%%%%%%%%%%%%
  \tikzifexternalizing{}{%
  \begin{tikzpicture}[remember picture,overlay]
  \node[anchor=south,yshift=-2mm] at (current page.south) 
  {\animategraphics[width=0.9995\paperwidth,autoplay,loop]{36}{assets/logo_TU_169_}{0}{39}};
    \end{tikzpicture}%
  }%
  % if you don't like the animation or it doesn't work for you, use this here instead:
%   \includegraphics[width=0.9995\paperwidth]{assets/logo_TU_169_0.pdf}
  %%%%%%%%%%%%%% END OF ANIMATED LOGO %%%%%%%%%%%

\end{frame}
\tikzexternalenable

\setlength{\figwidth}{.9\textwidth}
\setlength{\figheight}{.6\textheight}

%  \begin{frame}\frametitle{title}
%      \framesubtitle{subtitle}

%  \begin{itemize}
%   \item items
%   \item more items
%  \end{itemize}

%  \begin{block}{Block title}
%  a block
%  \end{block}

%  \ribbon{a ribbon across the page (for big takeaways).}

%  \end{frame}

%  \blackslidetext{This is a black slide for summaries or take-aways while the next slide is to draw attention to the speaker since there is no visual input for the listeners}
%  %\blackslide

\begin{frame}{Introduction}
\framesubtitlecite{Tumor segmentation}{\cite{AlBadawy2018}}
\begin{figure}
    \centering
    \includegraphics[width=\textwidth]{fig/medical.PNG}
\end{figure}
\end{frame}

\begin{frame}{Introduction}
\framesubtitlecite{Tumor segmentation}{\cite{AlBadawy2018}}
\begin{figure}
    \centering
    \includegraphics[width=0.85\textwidth]{fig/tumor.PNG}
\end{figure}
\end{frame}

\begin{frame}{Introduction}
\framesubtitle{High-Level Task Overview}
\begin{figure}
    \centering
    \includegraphics[width=0.65\textwidth]{fig/DG.pdf}
\end{figure}
\end{frame}

\begin{frame}{Introduction}
\framesubtitlecite{Learning setups}{\cite{gulrajani2020search}}
    \begin{table}
    \centering
    \begin{tabular}{lll}
    \toprule
    \textbf{Setup} & \textbf{Training inputs}  & \textbf{Testing inputs} \\
    \midrule
        Generative learning & $U_{\env_1}$ & $\emptyset$ \\ 
        Unsupervised learning & $U_{\env_1}$ & $U_{\env_1}$  \\ 
        Supervised learning & $L_{\env_1}$ & $U_{\env_1}$ \\ 
        Semi-supervised learning & $L_{\env_1}, U_{\env_1}$ & $U_{\env_1}$ \\ 
        Multitask learning & $L_{\env_1}, \ldots, L_{\env_s}$ & $U_{\env_1}, \ldots, U_{\env_s}$ \\ 
        Continual (lifelong) learning & $L_{\env_1}, \ldots, L_{\env_\infty}$ & $U_{\env_1}, \ldots, U_{\env_\infty}$ \\ 
        Domain Adaptation & $L_{\env_1}, \ldots, L_{\env_s}, U_{\env_t}$ & $U_{\env_t}$ \\ 
        Transfer learning & $U_{\env_1}, \ldots, U_{\env_s}, L_{\env_t}$ & $U_{\env_t}$ \\ 
        \emph{Domain generalization} & $L_{\env_1}, \ldots, L_{\env_s}$ & $U_{\env_t}$ \\ 
    \bottomrule
    \end{tabular}
    \caption[Differences in learning setups]{Differences in learning setups adapted from: \cite{gulrajani2020search}. For each environment $\env$ the labeled and unlabeled datasets are denoted as $L_\env$ and $U_\env$ respectively.}
    \label{tab:learning_setups}
\end{table}
\end{frame}

\begin{frame}{Introduction}
\framesubtitlecite{Samples across popular datasets}{}
\begin{table}
    \centering
    \begin{tabularx}{\textwidth}{lcYYYY}
    \toprule
    \textbf{Dataset}  & \textbf{Reference} & $\env^1$   &  $\env^2$ & $\env^3$ & $\env^4$ \\
    \midrule
          %%
       \multirow{7}{*}{\textbf{Office-Home}} & \multirow{7}{*}{\cite{VenkateswaraECP17}} & \domainsize{Art} & \domainsize{Clipart} & \domainsize{Product} & \domainsize{Real} \\
       & & \includegraphics[height=\imagequadsize, width=\imagequadsize]{fig/OfficeHome_env0Art_1_idx61_class0.png} & \includegraphics[height=\imagequadsize, width=\imagequadsize]{fig/OfficeHome_env1Clipart_5_idx24_class0.png} & \includegraphics[height=\imagequadsize, width=\imagequadsize]{fig/OfficeHome_env2Product_14_idx67_class0.png} & \includegraphics[height=\imagequadsize, width=\imagequadsize]{fig/OfficeHome_env3Real World_44_idx13_class0.png} \\
        & & \includegraphics[height=\imagequadsize, width=\imagequadsize]{fig/OfficeHome_env0Art_18_idx176_class3.png} & \includegraphics[height=\imagequadsize, width=\imagequadsize]{fig/OfficeHome_env1Clipart_12_idx199_class3.png} & \includegraphics[height=\imagequadsize, width=\imagequadsize]{fig/OfficeHome_env2Product_19_idx261_class3.png} & \includegraphics[height=\imagequadsize, width=\imagequadsize]{fig/OfficeHome_env3Real World_20_idx300_class3.png} \\
        \addlinespace
        %%
       \multirow{7}{*}{\textbf{VLCS}} & \multirow{7}{*}{\cite{FangXR13}} & \domainsize{V} & \domainsize{L} & \domainsize{C} & \domainsize{S} \\ 
       & & \includegraphics[height=\imagequadsize, width=\imagequadsize]{fig/voc_bird.jpg}  &  \includegraphics[height=\imagequadsize, width=\imagequadsize]{fig/labelme_bird.jpg} & \includegraphics[height=\imagequadsize, width=\imagequadsize]{fig/caltech_bird.jpg} & \includegraphics[height=\imagequadsize, width=\imagequadsize]{fig/sun_bird.jpg} \\
       & & \includegraphics[height=\imagequadsize, width=\imagequadsize]{fig/voc_car.jpg} & \includegraphics[height=\imagequadsize, width=\imagequadsize]{fig/labelme_car.jpg} & \includegraphics[height=\imagequadsize, width=\imagequadsize]{fig/caltech_car.jpg} & \includegraphics[height=\imagequadsize, width=\imagequadsize]{fig/sun_car.jpg} \\
    \bottomrule
    \end{tabularx}
    \caption{Samples for two different classes across domains for popular datasets}
    \label{tab:common_examples}
\end{table}
\end{frame}



\begin{frame}{Introduction}
\framesubtitlecite{Samples across popular datasets}{}
\begin{table}
    \centering
    \begin{tabularx}{\textwidth}{lcYYYY}
    \toprule
    \textbf{Dataset}  & \textbf{Reference} & $\env^1$   &  $\env^2$ & $\env^3$ & $\env^4$ \\
    \midrule
       %%
       \multirow{7}{*}{\textbf{Terra Incognita}} & \multirow{7}{*}{\cite{BeeryHP18}} & \domainsize{L100} & \domainsize{L38} & \domainsize{L43} & \domainsize{L46} \\ 
       & & \includegraphics[height=\imagequadsize, width=\imagequadsize]{fig/TerraIncognita_env0location_100_22_idx685_class4.png} &  \includegraphics[height=\imagequadsize, width=\imagequadsize]{fig/TerraIncognita_env1location_38_5_idx2429_class4.png} &
       \includegraphics[height=\imagequadsize, width=\imagequadsize]{fig/TerraIncognita_env2location_43_45_idx2058_class4.png} &  \includegraphics[height=\imagequadsize, width=\imagequadsize]{fig/TerraIncognita_env3location_46_6_idx2012_class4.png} \\
       & & \includegraphics[height=\imagequadsize, width=\imagequadsize]{fig/TerraIncognita_env0location_100_0_idx823_class5.png} & \includegraphics[height=\imagequadsize, width=\imagequadsize]{fig/TerraIncognita_env1location_38_35_idx2747_class5.png} & \includegraphics[height=\imagequadsize, width=\imagequadsize]{fig/TerraIncognita_env2location_43_38_idx2211_class5.png} & \includegraphics[height=\imagequadsize, width=\imagequadsize]{fig/TerraIncognita_env3location_46_37_idx2142_class5.png} \\
       \addlinespace
       %%
       \multirow{7}{*}{\textbf{PACS}} & \multirow{7}{*}{\cite{LiYSH17}}  & \domainsize{Art} & \domainsize{Cartoon} & \domainsize{Photo} & \domainsize{Sketch} \\
       & & \includegraphics[height=\imagequadsize, width=\imagequadsize]{fig/dog_1.jpg} &  \includegraphics[height=\imagequadsize, width=\imagequadsize]{fig/dog_2.jpg} &  \includegraphics[height=\imagequadsize, width=\imagequadsize]{fig/dog_3.jpg} &  \includegraphics[height=\imagequadsize, width=\imagequadsize]{fig/dog_4.png} \\
       & & \includegraphics[height=\imagequadsize, width=\imagequadsize]{fig/elephant_1.jpg} & \includegraphics[height=\imagequadsize, width=\imagequadsize]{fig/elephant_2.jpg} & \includegraphics[height=\imagequadsize, width=\imagequadsize]{fig/elephant_3.jpg} & \includegraphics[height=\imagequadsize, width=\imagequadsize]{fig/elephant_4.png} \\
    \bottomrule
    \end{tabularx}
    \caption{Samples for two different classes across domains for popular datasets}
    \label{tab:common_examples}
\end{table}
\end{frame}


\begin{frame}{Introduction}
\framesubtitlecite{The intuitive baseline}{}
\emph{Ordinary empirical risk:}
\begin{equation}
    \riske =\frac{1}{n} \sum_{i=1}^n \lossi{\f{\xxi}}{\yyi},
\end{equation}
\vspace{0.5cm}

\uncover<2->{\emph{Assumption:} Minimizing the empirical risk over all source domains achieves good generalization
\begin{equation}
\label{eq:domain_risk_emp}
    \riske = \frac{1}{s} \sum_{\env \in \envs} \frac{1}{n_\env} \sum_{i=1}^{n_\env} \lossi{\f{\xxi^{\env}}} {\yyi^{\env}}
\end{equation}
}
\end{frame}

\begin{frame}{Related Work}
\framesubtitle{What have other people tried?}
    \begin{itemize}
        \item \emph{Learning invariant features:} Kernel methods, adversarial approaches
        \item \emph{Model Ensembling:} Domain-specific batch norm, predictors
        \item \emph{Meta-Learning:} Learning the domain-shift
        \item \emph{Data Augmentation:} \textsc{Mixup} merge samples, remove textural information
    \end{itemize}
\end{frame}

\begin{frame}{Related Work}
\framesubtitlecite{What have other people tried?: Model-Ensembling}{\cite{ManciniBC018}}
\includegraphics[width=\textwidth]{fig/title_dg.PNG}
\end{frame}

\begin{frame}{Related Work}
\framesubtitlecite{What have other people tried?: Model-Ensembling}{\cite{ManciniBC018}}
\begin{figure}
    \centering
    \includegraphics[width=0.7\textwidth]{fig/method-dg2.pdf}
    \caption{Ensemble Method Overview}
    \label{fig:my_label}
\end{figure}
\end{frame}

\begin{frame}{Related Work}
\framesubtitlecite{What have other people tried?: Data Augmentation}{\cite{RahmanFBS19}}
\includegraphics[width=\textwidth]{fig/aug_dg.PNG}
\end{frame}

\begin{frame}{Related Work}
\framesubtitlecite{What have other people tried?: Data Augmentation}{\cite{RahmanFBS19}}
\begin{figure}
    \centering
    \includegraphics[width=0.49\textwidth]{fig/GAN_Images_f.pdf}
    \includegraphics[width=0.49\textwidth]{fig/WACV1_1_ff.pdf}
    \caption{Data Augmentation Overview}
    \label{fig:my_label}
\end{figure}
\end{frame}

\blackslidetext{\centering \large How does Representation Self Challenging approach this task?}

\begin{frame}{Representation Self-Challenging}
\framesubtitlecite{Relationship to Explainability}{\cite{huang2020selfchallenging}}
\centering
 \includegraphics[width=0.5\textwidth]{fig/RSCExplainability_1.pdf}
\end{frame}

\begin{frame}{Representation Self-Challenging}
\framesubtitlecite{Relationship to Explainability}{\cite{huang2020selfchallenging}}
\centering
 \includegraphics[width=0.5\textwidth]{fig/RSCExplainability_2.pdf}
\end{frame}

\begin{frame}{Representation Self-Challenging}
\framesubtitlecite{Relationship to Explainability}{\cite{huang2020selfchallenging}}
\centering
 \includegraphics[width=0.5\textwidth]{fig/RSCExplainability_3.pdf}
\end{frame}

\begin{frame}{Representation Self-Challenging}
\framesubtitlecite{Relationship to Explainability}{\cite{huang2020selfchallenging}}
\centering
 \includegraphics[width=0.5\textwidth]{fig/RSCExplainability_4.pdf}
\end{frame}

\begin{frame}{Representation Self-Challenging}
\framesubtitlecite{A high level overview}{\cite{huang2020selfchallenging}}
    \begin{figure}
        \centering
        \includegraphics[width=0.65\textwidth]{fig/intro.pdf}
        \caption{Overview of Representation Self-Challenging}
        \label{fig:my_label}
    \end{figure}
\end{frame}

\begin{frame}{Representation Self-Challenging}
\framesubtitlecite{What they actually do}{}
    \begin{figure}
        \centering
        \includegraphics[width=\textwidth]{fig/rsc_model_spatial_avg.pdf}
        \caption{Overview of Representation Self-Challenging: Spatial-Wise Average}
    \end{figure}
\end{frame}

\begin{frame}{Representation Self-Challenging}
\framesubtitlecite{What they actually do}{}
    \begin{figure}
        \centering
        \includegraphics[width=\textwidth]{fig/rsc_model_channel_avg.pdf}
        \caption{Overview of Representation Self-Challenging: Channel-Wise Average}
    \end{figure}
\end{frame}

\begin{frame}{Representation Self-Challenging}
\framesubtitlecite{What they actually do}{}
    \begin{figure}
        \centering
        \includegraphics[width=\textwidth]{fig/rsc_change.pdf}
        \caption{Overview of Representation Self-Challenging: Mask Permutation}
    \end{figure}
\end{frame}

\begin{frame}{Representation Self-Challenging}
\framesubtitlecite{What they actually do}{}
\begin{algorithm}[H]
\small
    \SetAlgoLined
    \SetNoFillComment
    \SetKwInOut{Input}{Input}
    \Input{Data $\mathbf{X}, \mathbf{Y}$ with $\mathbf{x}_i \in \mathbb{R}^{H \times W \times 3}$, drop factor $f$, \textcolor{MPG}{$b$}, epochs $T$}
    %\KwResult{how to write algorithm with \LaTeX2e }
    %initialization\;
    \BlankLine
    \While{$epoch \leq T$}{
    \For{every sample (or batch) $\mathbf{x}, \mathbf{y}$}{
    Extract features $\mathbf{z} = \phi(\mathbf{x})$ \tcp*[r]{$\mathbf{z}$ has shape  $\mathbb{R}^{H_\mathbf{z} \times W_\mathbf{z} \times K} $}
    Compute gradient $\featureg$ w.r.t features  \;
    \textcolor{MPG}{Compute $\featuregal$ by avg. pooling using two different methods}\;
    \textcolor{MPG}{Duplicate $\featurega$ along channel/spatial dimension for initial shape}\;
    Compute mask $\mathbf{m}_{i,j}$ \;
    Mask features to obtain $\Tilde{\mathbf{z}} = \mathbf{m} \odot \mathbf{z}$ \tcp*[r]{Evaluate effect of preliminary mask $\downarrow$}
    \textcolor{MPG}{Compute change $\mathbf{c}$} \;
    \textcolor{MPG}{Revert masking for specific samples} \;
    Mask features $\Tilde{\mathbf{z}} = \mathbf{z} \odot \mathbf{m}$ \;
    Compute loss $\mathcal{L}(w(\Tilde{\mathbf{z}}), \mathbf{y})$ and backpropagate to whole network
    }
}
\caption{Spatial- and Channel-Wise RSC}
\label{alg:SpatialRSC}
\end{algorithm}
\end{frame}


\blackslidetext{\centering \large Is this the best we can do in terms of pushing the network to different discriminative features? What about using explainability methods to guide the network?
}

\begin{frame}{Diversified Class Activation Maps (\divcam)}
\framesubtitlecite{Our first approach}{}
    \begin{figure}
        \centering
        \includegraphics[width=\textwidth]{Figures/Chapter4/model_figure-cropped.pdf}
        \caption{Overview of Diversified Class Activation Maps}
    \end{figure}
    \vspace{0.2cm}
    \uncover<2->{\ribbon{By reconstructing to Grad-CAMs we properly take into account the content of the image!}}
\end{frame}

\begin{frame}{Results}
\framesubtitlecite{Outside of \domainbed}{}
    \begin{table}
\small
\centering
\begin{tabular}{llcccccc}
\toprule
\textbf{Algorithm} & \textbf{Ref.} & \textbf{Backbone} & \textbf{P} & \textbf{A} & \textbf{C} & \textbf{S} &  \textbf{Avg.} \\
\midrule
\textsc{Baseline}		&\cite{CarlucciDBCT19}				&	ResNet-18	&	$95.73$		&	$77.85$		&	$74.86$		&	$67.74$		&	$79.05$		 \\
\textsc{MASF}		&\cite{DouCKG19}					&	ResNet-18	&	$94.99$		&	$80.29$		&	$77.17$		&	$71.69$		&	$81.03$		 \\
\textsc{Epi-FCR}		&\cite{LiZYLSH19}					&	ResNet-18	&	$93.90$		&	$82.10$		&	$77.00$		&	$73.00$		&	$81.50$		 \\
\textsc{JiGen}		&\cite{CarlucciDBCT19}				&	ResNet-18	&	$96.03$		&	$79.42$		&	$75.25$		&	$71.35$		&	$80.51$		\\
\textsc{MetaReg}		& \cite{BalajiSC18}					&	ResNet-18	&	$95.50$		&	$\tabtop{83.70}$		&	$77.20$		&	$70.30$		&	$81.70$		\\
\textsc{RSC} (reported)	& \cite{huang2020selfchallenging}		&	ResNet-18	&	$95.99$		&	$83.43$		&	$80.31$		&	$80.85$		&	$85.15$		\\
\textsc{RSC} (reproduced) & \cite{huang2020selfchallenging}		&	ResNet-18	&	$93.73$		&	$80.41$		&	$77.53$		&	$80.79$		&	$83.12$		\\
\divcams			&     (ours)						&	ResNet-18	&	$\tabtop{96.11}$		&	$80.27$		&	$\tabtop{77.82}$		&	$\tabtop{82.18}$		&	$\tabtop{84.10}$		\\
\bottomrule
\end{tabular}
\caption{Performance comparison for PACS outside of the \domainbed framework with the official data split.}
\label{tab:official}
\end{table}
\end{frame}

\begin{frame}{Results}
\framesubtitlecite{Within \domainbed}{}
\begin{figure}
    \centering
    \includegraphics[width=0.8\textwidth]{Presentation/fig/results.PNG}
    \caption{Performance comparison for PACS within the \domainbed framework.}
    \label{fig:my_label}
\end{figure}
\end{frame}

\begin{frame}{Class Activation Maps}
\framesubtitlecite{Visualizations of the CAMs throughout training}{}
    \begin{figure}
    \centering
    \begin{tabularx}{\textwidth}{YYYY}
       \textbf{Original} & \textbf{Step 300} & \textbf{Step 2700} & \textbf{Step 4500} \\[0.2cm]
         \includegraphics[width=\imagequadsizecams]{Figures/Chapter4/img66.png} &  \includegraphics[width=\imagequadsizecams]{Figures/Chapter4/giraffe1_step_300_cam.png} & \includegraphics[width=\imagequadsizecams]{Figures/Chapter4/giraffe1_step_2700_cam.png} &
        \includegraphics[width=\imagequadsizecams]{Figures/Chapter4/giraffe1_step_4500_cam.png} \\
       \includegraphics[width=\imagequadsizecams]{Figures/Chapter4/img76.png} &
        \includegraphics[width=\imagequadsizecams]{Figures/Chapter4/elephant1_step_300_cam.png} &
        \includegraphics[width=\imagequadsizecams]{Figures/Chapter4/elephant1_step_2700_cam.png} &
        \includegraphics[width=\imagequadsizecams]{Figures/Chapter4/elephant1_step_4500_cam.png} \\
    \end{tabularx}
    \caption[Used class activation maps in \divcams throughout training]{Used class activation maps for \divcams in update step $300/5000$, $2700/5000$, and $4500/5000$ using a ResNet-50 backbone.}
    \label{fig:cams_and_masks_divcam}
\end{figure}
\end{frame}

\begin{frame}{Some improvements for the class activation maps}
\framesubtitlecite{Global Average Pooling bias: Threshold average pooling layer}{\cite{Bae2020RethinkingCAM}}
    \begin{figure}
        \centering
        \includegraphics[width=\textwidth]{Presentation/fig/small_activations.PNG}
        \caption{Problems of the GAP layer: \cite{Bae2020RethinkingCAM}}
    \end{figure}
     \uncover<2->{\ribbon{They propose an average pooling which is a tradeoff between max and avg pooling}}
\end{frame}

\begin{frame}{Some improvements for the class activation maps}
\framesubtitlecite{Global Average Pooling bias: Smoothing negative CAMs}{\cite{sun2020fixing}}
    \begin{figure}
        \centering
        \includegraphics[width=0.6\textwidth]{Presentation/fig/neg_maps.PNG}
        \caption{Problems of the negative CAMs: \cite{sun2020fixing}}
    \end{figure}
     \uncover<2->{\ribbon{They propose minimizing the KL-divergence between an uniform map and the negative CAMs}}
\end{frame}

\begin{frame}{Ablation study for the ``improvements''}
\framesubtitlecite{These approaches do not help for training-domain validation}{}
    \begin{table}
    \footnotesize
%\renewcommand{\arraystretch}{1.2}
    \centering
    \begin{tabular}{lccccc}
    \toprule
    \textbf{Name} &  \textbf{P} & \textbf{A} & \textbf{C} & \textbf{S} & \textbf{Avg.} \\
    \midrule 
    \divcam & $\tabtop{97.6\pm0.4}$ & $85.2\pm0.8$ & $80.5\pm0.7$ & $78.3\pm0.8$ & $\tabtop{85.4\pm0.5}$ \\
    \divcams & $97.3\pm0.4$ & $86.2\pm1.4$ & $79.1\pm2.2$ & $\tabtop{79.2\pm0.1}$ & $85.4\pm0.2$ \\
    \divcams + TAP & $96.9\pm0.1$ & $85.1\pm1.5$ & $78.7\pm0.4$ & $75.3\pm0.6$ & $84.0\pm0.4$ \\
    \divcams + HNC & $97.2\pm0.3$ & $\tabtop{87.2\pm0.9}$ & $79.2\pm0.6$ & $71.7\pm3.1$ & $83.8\pm0.4$ \\
    \divcams + CDANN & $97.5\pm0.4$ & $85.2\pm2.8$ & $78.3\pm2.0$ & $74.8\pm0.9$ & $84.0\pm1.5$ \\
    \divcams + MMD & $97.0\pm0.2$ & $85.4\pm1.0$ & $\tabtop{81.5\pm0.4}$ & $75.8\pm3.5$ & $84.9\pm1.1$ \\
    CAM + CDANN & $97.2\pm0.3$ & $86.7\pm0.5$ & $77.3\pm1.7$ & $71.5\pm1.3$ & $83.2\pm0.8$ \\
    \midrule
    \tdivcam & $96.2\pm1.2$ & $87.0\pm0.5$ & $82.0\pm0.9$ & $80.8\pm0.6$ & $86.5\pm0.1$ \\
    \tdivcams & $97.2\pm0.3$ & $86.5\pm0.4$ & $83.0\pm0.5$ & $82.2\pm0.1$ & $87.2\pm0.1$ \\
    \divcams + TAP* & $97.3\pm0.3$ & $87.2\pm0.8$ & $\tabtop{83.2\pm0.8}$ & $\tabtop{82.8\pm0.2}$ & $\tabtop{87.6\pm0.0}$ \\
    \divcams + HNC*  & $97.3\pm0.2$ & $\tabtop{87.4\pm0.5}$ & $81.4\pm0.6$ & $79.7\pm1.1$ & $86.5\pm0.4$ \\
    \divcams + CDANN* & $\tabtop{97.3\pm0.5}$ & $85.9\pm1.2$ & $80.6\pm0.4$ & $80.9\pm0.4$ & $86.2\pm0.2$ \\
    \divcams + MMD* & $\tabtop{97.3\pm0.5}$ & $86.8\pm0.7$ & $83.2\pm0.4$ & $80.9\pm0.7$ & $87.1\pm0.4$ \\
    CAM + CDANN* & $97.2\pm0.4$ & $86.7\pm0.5$ & $81.9\pm0.2$ & $80.6\pm0.7$ & $86.6\pm0.2$ \\
    \bottomrule
    \end{tabular}
    \caption{Ablation study for the \divcam masks on the PACS dataset using training-domain validation (top) and oracle validation denoted with * (bottom). We use a ResNet-50 backbone and $20$ hyperparameter samples.}
    \label{tab:scam_masks}
\end{table}
\end{frame}

\blackslidetext{\centering \large How about applying prototypical networks to domain generalization?}

\begin{frame}{Prototypical Networks for Domain Generalization}
\framesubtitlecite{An intuitive baseline}{\cite{ChenLTBRS19}}
    \begin{figure}
    \centering
    \includegraphics[width=\textwidth]{Figures/Chapter4/prototype_agnostic-cropped.pdf}
    \caption{Domain-agnostic Prototype Layer}
    \label{fig:my_label}
\end{figure}
\end{frame}

\begin{frame}{Prototypical Networks for Domain Generalization}
\framesubtitlecite{An intuitive baseline}{\cite{ChenLTBRS19}}
Domain-Agnostic cluster and separation losses:
   \begin{alignat}{3}
\label{eq:prototype_losses}
    \mathcal{L}_{\mathrm{clst}} &= &&\frac{1}{n} \sum_{i=1}^n \min_{j: \prot \in \prots_{\yyi}} \min_{\zpatch \in \text{patches}(\zz)} \left\|\zpatch - \prot\right\|^2_2\\
    \mathcal{L}_{\mathrm{sep}} &= -&&\frac{1}{n} \sum_{i=1}^n \min_{j: \prot \notin \prots_{\yyi}} \min_{\zpatch \in \text{patches}(\zz)} \left\|\zpatch - \prot\right\|^2_2.
\end{alignat}
\end{frame}

\begin{frame}{Prototypical Networks for Domain Generalization}
\framesubtitlecite{Sometimes not all prototypes are used properly}{}
    \begin{figure}
    \centering
    \includegraphics[width=0.9\textwidth]{Figures/Chapter4/2021-01-21-ProDropIncorrectWeight-1.0SAVEResNet18oracle_validation_trial1.pdf}
    \caption{Pairwise learned prototype distances}
    \label{fig:my_label}
\end{figure}
\end{frame}

\begin{frame}{Improvements: Prototype Dropping (\prodrop)}
\framesubtitlecite{Sometimes not all prototypes are used properly}{}
    \begin{algorithm}[H]
    \SetAlgoLined
    \SetKwInOut{Input}{Input}
    \Input{Data $\mathbf{X}, \mathbf{Y}$ with $\mathbf{x}_i \in \mathbb{R}^{H \times W \times 3}$, drop factor $p,b$, epochs $T$}
    \BlankLine
    \While{$epoch \leq T$}{
        \For{every batch $\mathbf{x}, \mathbf{y}$}{
            Extract features $\mathbf{z} = \phi(\mathbf{x})$ \tcp*[r]{$\mathbf{z}$ has shape  $\mathbb{R}^{H_\mathbf{z} \times W_\mathbf{z} \times K} $}
            Compute $\unit(\zz)$\;
            Compute $\mathbf{m}_{c,j}$ \;
            Adapt $\mathbf{m}_{c,j}$ \;
            Compute $\mplayer(\zz)$ \;
            Backpropagate loss $\mathcal{L}_{ce}(w(\mplayer(\zz)), \mathbf{y}) + \lambda_4 \mathcal{L}_{\mathrm{clst}} + \lambda_5 \mathcal{L}_{\mathrm{sep}}$ \;
            }
    }
\caption{Prototype Dropping (\prodrop)}
\label{alg:ProDrop}
\end{algorithm}
\end{frame}

\begin{frame}{Prototypical Networks for Domain Generalization}
\framesubtitlecite{Sometimes not all prototypes are used properly}{}
    \begin{figure}
    \centering
    \includegraphics[width=0.9\textwidth]{Figures/Chapter4/2021-01-21-ProDropIncorrectWeight-1.0WithSCdrop_f0.5SAVEResNet18oracle_validation_trial1.pdf}
    \caption{Pairwise learned prototype distances}
    \label{fig:my_label}
\end{figure}
\end{frame}

\begin{frame}{Improvements: Intra-Loss}
\framesubtitlecite{Sometimes not all prototypes are used properly}{}
\begin{equation}
\label{eq:intra_loss}
    \mathcal{L}_\mathrm{intra} = \sum_{\proti_i, \prot \in \prots_c} \lambda_{\ell_2} \underbrace{\left\|\proti_i - \prot  \right\|_2 \vphantom{\frac{\proti_i\prot}{\left\|\proti_i \right\|_2 \left\|\prot \right\|_2}}}_{\ell_2-\text{distance}} + \lambda_{\cdistance} \underbrace{ (1-\frac{\proti_i\prot}{\left\|\proti_i \right\|_2 \left\|\prot \right\|_2})}_{\text{cosine-distance}},
\end{equation}
\end{frame}

\begin{frame}{Results on \prodrop and Intra-Loss}
\framesubtitlecite{Self-Challenging helps!}{}
    \begin{figure}
        \centering
        \includegraphics[width=0.7\textwidth]{Presentation/fig/intra_ablat.PNG}
        \label{fig:my_label}
    \end{figure}
\end{frame}

\begin{frame}{Prototypical Networks for Domain Generalization}
\framesubtitlecite{Incorporating domain information}{}
    \begin{figure}
    \centering
    \includegraphics[width=\textwidth]{Figures/Chapter4/prototype_ensemble-cropped.pdf}
    \caption{Ensemble Prototype Layer}
    \label{fig:my_label}
\end{figure}
\end{frame}

\blackslidetext{\centering \large So far the prototypes have been directly learned, but we can also utilize a support set of images!}


\begin{frame}{Prototypical Networks for Domain Generalization}
\framesubtitlecite{Support Set Prototypes with attention aligment: crossTransformers}{\cite{DoerschGZ20}}
    \begin{figure}
    \centering
    \includegraphics[width=0.8\textwidth]{Presentation/fig/ctx-arch.pdf}
    \caption{Support Set Prototypes with attention aligment}
    \label{fig:my_label}
\end{figure}
\ribbon{We can do this for each domain separately and combine the final distance values!}
\end{frame}

\begin{frame}{Results}
\framesubtitlecite{\domainbed results for \dtransformers}{}
    \begin{figure}
    \centering
    \includegraphics[width=0.8\textwidth]{Presentation/fig/res_full.PNG}
    \caption{Full results for training-domain validation within \domainbed}
    \label{fig:my_label}
\end{figure}
\end{frame}

\begin{frame}{Results}
\framesubtitlecite{\domainbed results for \dtransformers}{}
    \begin{figure}
    \centering
    \includegraphics[width=0.8\textwidth]{Presentation/fig/res_full_2.PNG}
    \caption{Full results for training-domain validation within \domainbed}
    \label{fig:my_label}
\end{figure}
\end{frame}

\begin{frame}{Results}
\framesubtitlecite{\domainbed results for \dtransformers}{}
    \begin{figure}
    \centering
    \includegraphics[width=0.8\textwidth]{Presentation/fig/res_full_test.PNG}
    \caption{Full results for oracle-domain validation within \domainbed}
    \label{fig:my_label}
\end{figure}
\end{frame}

\begin{frame}{Results}
\framesubtitlecite{\domainbed results for \dtransformers}{}
    \begin{figure}
    \centering
    \includegraphics[width=0.8\textwidth]{Presentation/fig/res_full_test_2.PNG}
    \caption{Full results for oracle-domain validation within \domainbed}
    \label{fig:my_label}
\end{figure}
\end{frame}

\begin{frame}{Conclusion \& Outlook}
\framesubtitlecite{Some things to think about}{}
    \begin{itemize}
        \item \divcam works well and outperforms \rsc
        \uncover<2->{\item \prodrop and \dtransformers can achieve very good performance depending on the backbone and dataset}
        \uncover<3->{\item Upscaling $7 \times 7$ feature map to $14 \times 14$ to achieve big performance gains}
        \uncover<4->{\item Even though \prodrop does well on ResNet-18 it doesn't generalize well to ResNet-50}
        \uncover<5->{\item Clear ablation of the distance functions (maybe also metric learning)}
        \uncover<6->{\item More sophisticated method for aggregating across domains in \dtransformers}
    \end{itemize}
    \vspace{1cm}
    \uncover<7->{\ribbon{Overall: First stepping stone and analysis on explainability methods for DG, hopefully more to come!}}
\end{frame}


\begin{frame}[allowframebreaks]{References}
    \bibliographystyle{apalike}
    \bibliography{thesis}
\end{frame}

\end{document}